\section{Introduction}\label{sec:intro}

Public metabolomics deposits support method development, benchmarking and model evaluation~\citep{haug2020metabolights,sud2016workbench}. These uses require an analyst to identify which samples share a participant. Repeated measurements can support validation if dependence is modelled or one prespecified observation is selected per participant. Treating rows as independent without establishing their sampling units risks pseudoreplication~\citep{hurlbert1984pseudoreplication,lazic2010pseudoreplication}. Likewise, deposits that share participants cannot serve as independent validation units~\citep{riley2016external}.

Transcriptomic studies have documented hidden duplicates that can inflate cross-study validation and widespread sample misannotation~\citep{waldron2016doppelganger,toker2016whose}. Controlled-access genomic archives address linkage directly: dbGaP requires a subject--sample mapping and permits several samples per subject~\citep{dbgap_submission}. Metabolomics audits have found uneven compliance with reporting standards, variable metadata quality, mismatched sample identifiers and limited controlled-vocabulary coverage~\citep{spicer2017compliance,spicer2017decade,ferreira2017assessing,powell2021mwtab,jarmusch2020redu}. We are unaware of a metabolomics audit that tests whether available documentation identifies independent sampling units under a prespecified evidence rule.

We apply an affirmative evidence standard in four analyses with distinct denominators: a metadata census of 3,582 public MetaboLights accessions; a factor-metadata audit of 104 selected human Metabolomics Workbench studies; a documentary assessment of 40 randomly selected unresolved MetaboLights candidates; and a pre-verification of 19 human MetaboLights deposits excluded by a quality-control rule. The findings measure what the available documentation establishes about sampling.
